\documentclass[11pt]{scrartcl}
\usepackage[sexy]{evanptgs}
\usepackage{braket}
\numberwithin{example}{section}

\begin{document}

\title{Algoritmo de Shor: Criptografia Quântica}
\author{Marcos Burdzinski, Pedro Duailibe}
\date{}
\maketitle
\newpage
\tableofcontents
\newpage

\section{Introdução e Notação}

\subsection{Estados quânticos e medições}

Aqui utilizaremos a notação \textbf{bra-ket} (muito utilizada por físicos), onde um vetor é denotado por $\ket \bullet$, onde $\bullet$ é um nome de variável. 

\begin{example}
   Se $H$ é um espaço de Hilbert, com $\dim H = n < \infty$, os elementos base de $H$ são comumente definidos como
   \begin{align*}
     \ket 0, \ket 1, \dots , \ket {n-1}
   \end{align*}
   (em vez de $e_i$) e um elemento qualquer de $H$ é denotado por 
   \begin{align*}
       \ket \psi, \ket \phi, \dots
   \end{align*}
   
\end{example}

Para qualquer $\ket \psi \in H$, existe o seu dual canônico em $H ^\vee$, que denotamos por $\bra \psi$.

\begin{example}
   Se $\dim H = 2$, temos $\ket \phi = \begin{bmatrix} \alpha \\ \beta \end{bmatrix}$ em uma base ortonormal, então $\bra \psi = \begin{bmatrix} \overline \alpha & \overline \beta \end{bmatrix}$.

   Se $\ket \phi = \begin{bmatrix} \gamma \\ \delta \end{bmatrix}$, temos que o produto interno entre $\bra \psi$ e $\ket \phi$ é
   \begin{align*}
     \braket{\psi|\phi} = \begin{bmatrix} \overline \alpha & \overline \beta \end{bmatrix} \begin{bmatrix} \gamma \\ \delta \end{bmatrix} = \overline \alpha \gamma + \overline \beta \delta.
   \end{align*}
\end{example}

\subsubsection{O Espaço de Estados}

\begin{definition}
  Dado um espaço de Hilbert $H$ com $\dim H = n$, um \vocab{qubit} é um elemento não zero $\ket \psi \in H$.   
\end{definition}

Ou seja, se $H = \CC^{\oplus 2}$, um qubit é algo da forma
\begin{align*}
    \ket \psi = \alpha \ket 0 + \beta \ket 1,
\end{align*}
onde $\alpha$ e $\beta$ não podem ser ambos zero. 

\subsubsection{Observações}

A forma que observações funcionam é estranha e (como não somos físicos) iremos nos restringir somente a descrição matemática dessa \textit{quantum framework}:

\begin{itemize}
  \ii Escolha um operador Hermitiano $T:H \to H$, que é denotado de \vocab{observável}.
  \ii Considere os seus autovalores $\lambda_0, \dots, \lambda_{n-1}$, e autovetores correspondentes $\ket 0_T, \dots,$ $\ket {n-1}_T$. Podemos assumir que $\ket 0_T,\dots,\ket{n-1}_T$ formam uma base ortonormal de $H$.
  \ii Escreva 
  \begin{align*}
    \ket \psi = c_0 \ket 0_T + c_1 \ket 1_T + \cdots + c_{n-1}\ket{n-1}_T.
  \end{align*}
  \ii Então, a probabilidade de observar $\lambda_i$ é 
  \begin{align*}
    \frac {|c_i|^2}{|c_0|^2 + \cdots + |c_{n-1}|^2}.
  \end{align*}
  Isso é chamado de fazer uma \vocab{observação sobre} $T$.
\end{itemize}

\subsubsection{Entreleçamento}

\subsection{Circuitos Quânticos}

\subsubsection{Portas lógicas quânticas}

Em mecânica quântica, como podemos ter \textit{combinações lineares} de elementos base, nossas portas lógicas consistiram de \textit{transformações lineares}. 

\begin{definition}[Porta Lógica quântica]
    Portas lógicas quânticas são matrizes unitárias. 
\end{definition}

\end{document}
